\section{Appendix}

The complete source code is organized below by architectural layer. All source files are available in the project GitHub repository.

\textbf{GitHub Repository:} \url{https://github.com/mcittkmims/es-labs}

\subsection{Application Layer}

\subsubsection{Lab Selector Entry Point}

\begin{lstlisting}[language=C++, caption=src/main.cpp --- Lab selector (Lab 4 section), label=lst:app_main]
#elif defined(LAB4)
    #include "lab4_main.h"
// In setup():
    lab4Setup();
// In loop():
    lab4Loop();
\end{lstlisting}

\subsubsection{Lab 4 Entry Point}

\begin{lstlisting}[language=C++, caption=lab4\_main.h --- Lab 4 entry point interface, label=lst:lab4_main_h]
#ifndef LAB4_MAIN_H
#define LAB4_MAIN_H
void lab4Setup();
void lab4Loop();
#endif
\end{lstlisting}

\begin{lstlisting}[language=C++, caption=lab4\_main.cpp --- Lab 4 entry point implementation, label=lst:lab4_main_cpp]
#include "lab4_main.h"
#include "lab4_config.h"
#include "shared_state.h"
#include "task_input.h"
#include "task_control.h"
#include "task_display.h"
#include <Arduino.h>
#include <Arduino_FreeRTOS.h>
#include <stdio.h>
#include "StdioSerial.h"

void lab4Setup() {
    stdioSerialInit(9600);
    // Print startup banner (omitted for brevity)
    sharedStateInit();
    xTaskCreate(vTaskInput, "Input",
        TASK_INPUT_STACK, NULL, TASK_INPUT_PRIORITY, NULL);
    xTaskCreate(vTaskControl, "Control",
        TASK_CONTROL_STACK, NULL, TASK_CONTROL_PRIORITY, NULL);
    xTaskCreate(vTaskDisplay, "Display",
        TASK_DISPLAY_STACK, NULL, TASK_DISPLAY_PRIORITY, NULL);
}

void lab4Loop() { /* FreeRTOS tasks handle all logic */ }
\end{lstlisting}

\subsubsection{Configuration and Pin Mapping}

\begin{lstlisting}[language=C++, caption=lab4\_config.h --- Pin mapping and configuration, label=lst:lab4_config]
#ifndef LAB4_CONFIG_H
#define LAB4_CONFIG_H
#include <Arduino.h>
// Pin assignments
static const uint8_t PIN_RELAY     = 7;
static const uint8_t PIN_PWM_ACT   = 3;
static const uint8_t PIN_LED_GREEN = 8;
static const uint8_t PIN_LED_RED   = 9;
static byte KEYPAD_ROW_PINS[4] = {22, 23, 24, 25};
static byte KEYPAD_COL_PINS[4] = {26, 27, 28, 29};
static const uint8_t LCD_I2C_ADDR = 0x27;
// Timing
static const uint16_t TASK_INPUT_PERIOD_MS   = 50;
static const uint16_t TASK_CONTROL_PERIOD_MS = 100;
static const uint16_t TASK_DISPLAY_PERIOD_MS = 500;
// FreeRTOS task config
static const uint16_t TASK_INPUT_STACK   = 256;
static const uint16_t TASK_CONTROL_STACK = 256;
static const uint16_t TASK_DISPLAY_STACK = 512;
static const uint8_t TASK_INPUT_PRIORITY   = 3;
static const uint8_t TASK_CONTROL_PRIORITY = 2;
static const uint8_t TASK_DISPLAY_PRIORITY = 1;
// Conditioning parameters
static const uint8_t  ACT_MEDIAN_WINDOW = 5;
static const float    ACT_EWMA_ALPHA    = 0.4f;
static const float    ACT_RAMP_STEP     = 5.0f;
static const float    OVERLOAD_THRESHOLD_HIGH = 90.0f;
static const float    OVERLOAD_THRESHOLD_LOW  = 80.0f;
#endif
\end{lstlisting}

\subsubsection{Shared State Module}

\begin{lstlisting}[language=C++, caption=shared\_state.h --- Shared state interface, label=lst:shared_state_h_app]
struct ActuatorState {
    bool relayCommandOn, relayActualOn;
    uint8_t relayDebounceCount;
    float pwmCommandPercent, pwmConditioned, pwmRamped;
    uint8_t pwmRawValue;
    bool overloadAlert, inputModeAnalog;
    char inputBuffer[4];
    uint8_t inputBufferLen;
};
void sharedStateInit();
void sharedStateLock();
void sharedStateUnlock();
ActuatorState* sharedStateGet();
\end{lstlisting}

\subsubsection{Input Task}

\begin{lstlisting}[language=C++, caption=task\_input.cpp --- Keypad input task, label=lst:task_input_cpp]
#include "task_input.h"
#include "shared_state.h"
#include "lab4_config.h"
#include "KeypadInput.h"
#include <stdio.h>
#include <stdlib.h>

static KeypadInput keypad(KEYPAD_ROW_PINS, KEYPAD_COL_PINS);

void vTaskInput(void *pvParameters) {
    keypad.init();
    TickType_t xLastWake = xTaskGetTickCount();
    const TickType_t xPeriod = pdMS_TO_TICKS(TASK_INPUT_PERIOD_MS);
    for (;;) {
        char key = keypad.getKey();
        if (key != 0) {
            sharedStateLock();
            ActuatorState *s = sharedStateGet();
            switch (key) {
                case 'A': s->relayCommandOn = !s->relayCommandOn; break;
                case 'B': s->inputModeAnalog = true;
                          s->inputBufferLen = 0; break;
                case 'C': s->relayCommandOn = false;
                          s->pwmCommandPercent = 0.0f; break;
                case '#': if (s->inputModeAnalog && s->inputBufferLen > 0) {
                    s->inputBuffer[s->inputBufferLen] = '\0';
                    int val = atoi(s->inputBuffer);
                    if (val > 100) val = 100;
                    s->pwmCommandPercent = (float)val;
                    s->inputModeAnalog = false;
                } break;
                case '*': s->inputBufferLen = 0;
                          s->inputModeAnalog = false; break;
                default:
                    if (key >= '0' && key <= '9' && s->inputModeAnalog)
                        if (s->inputBufferLen < 3)
                            s->inputBuffer[s->inputBufferLen++] = key;
                    break;
            }
            sharedStateUnlock();
        }
        vTaskDelayUntil(&xLastWake, xPeriod);
    }
}
\end{lstlisting}

\subsubsection{Control Task}

\begin{lstlisting}[language=C++, caption=task\_control.cpp --- Actuator control task, label=lst:task_control_cpp]
#include "task_control.h"
#include "shared_state.h"
#include "lab4_config.h"
#include "Relay.h"
#include "PwmActuator.h"
#include "ActuatorConditioner.h"
#include "ThresholdAlert.h"
#include "Led.h"

static Relay relay(PIN_RELAY, true);
static PwmActuator pwmAct(PIN_PWM_ACT);
static Led ledGreen(PIN_LED_GREEN);
static Led ledRed(PIN_LED_RED);
static ActuatorConditioner conditioner(
    ACT_MEDIAN_WINDOW, ACT_EWMA_ALPHA,
    ACT_MIN_CLAMP, ACT_MAX_CLAMP, ACT_RAMP_STEP);
static ThresholdAlert overloadAlert(
    OVERLOAD_THRESHOLD_HIGH, OVERLOAD_THRESHOLD_LOW,
    OVERLOAD_DEBOUNCE_COUNT);

void vTaskControl(void *pvParameters) {
    relay.init(); pwmAct.init();
    ledGreen.init(); ledRed.init();
    overloadAlert.init(); ledGreen.turnOn();
    TickType_t xLastWake = xTaskGetTickCount();
    for (;;) {
        sharedStateLock();
        ActuatorState *s = sharedStateGet();
        // Relay debounce + control
        relay.setState(s->relayCommandOn);
        s->relayActualOn = relay.isOn();
        // PWM conditioning pipeline
        float output = conditioner.process(s->pwmCommandPercent);
        pwmAct.setDuty(output);
        s->pwmRamped = output;
        s->pwmRawValue = pwmAct.getRawPwm();
        // Overload alert
        overloadAlert.update(output);
        s->overloadAlert = overloadAlert.isAlertActive();
        // LED indicators
        if (s->overloadAlert) { ledRed.turnOn(); ledGreen.turnOff(); }
        else { ledRed.turnOff(); ledGreen.turnOn(); }
        sharedStateUnlock();
        vTaskDelayUntil(&xLastWake,
            pdMS_TO_TICKS(TASK_CONTROL_PERIOD_MS));
    }
}
\end{lstlisting}

\subsubsection{Display Task}

\begin{lstlisting}[language=C++, caption=task\_display.cpp --- LCD and serial reporting task, label=lst:task_display_cpp]
#include "task_display.h"
#include "shared_state.h"
#include "lab4_config.h"
#include "LcdDisplay.h"
#include <stdio.h>
#include <stdlib.h>

static LcdDisplay lcd(LCD_I2C_ADDR, LCD_COLS, LCD_ROWS);

void vTaskDisplay(void *pvParameters) {
    lcd.init();
    TickType_t xLastWake = xTaskGetTickCount();
    uint8_t reportCounter = 0;
    for (;;) {
        sharedStateLock();
        ActuatorState *s = sharedStateGet();
        bool relayOn = s->relayActualOn;
        float pwmRamp = s->pwmRamped;
        bool alert = s->overloadAlert;
        sharedStateUnlock();
        // LCD update
        char line1[17], line2[17];
        snprintf(line1, 17, "Relay:%-3s PWM%3d%%",
            relayOn ? "ON" : "OFF", (int)(pwmRamp + 0.5f));
        snprintf(line2, 17, "Out:%3d%% %s",
            (int)(pwmRamp + 0.5f), alert ? "!ALERT" : "  OK  ");
        lcd.showTwoLines(line1, line2);
        // Serial report every 2s
        if (++reportCounter >= 4) {
            reportCounter = 0;
            printf("\r\n--- Status Report ---\r\n");
            printf("Relay: %s | PWM: %d%%\r\n",
                relayOn ? "ON" : "OFF", (int)(pwmRamp+0.5f));
            printf("Alert: %s\r\n",
                alert ? "OVERLOAD" : "Normal");
        }
        vTaskDelayUntil(&xLastWake,
            pdMS_TO_TICKS(TASK_DISPLAY_PERIOD_MS));
    }
}
\end{lstlisting}

\subsection{Service Layer}

\subsubsection{Actuator Conditioner}

\begin{lstlisting}[language=C++, caption=ActuatorConditioner.h --- Command conditioning interface, label=lst:act_cond_h_app]
#ifndef ACTUATOR_CONDITIONER_H
#define ACTUATOR_CONDITIONER_H
#include "SignalConditioner.h"
class ActuatorConditioner {
public:
    ActuatorConditioner(uint8_t medianWindow, float ewmaAlpha,
        float minClamp, float maxClamp, float maxRampStep);
    float process(float rawCommand);
    float getRampedOutput() const;
    float getConditionedTarget() const;
    const SignalConditioner& getSignalConditioner() const;
    void reset();
private:
    SignalConditioner _conditioner;
    float _maxRampStep, _rampedOutput, _conditionedTarget;
    bool _initialized;
};
#endif
\end{lstlisting}

\begin{lstlisting}[language=C++, caption=ActuatorConditioner.cpp --- Conditioning with ramping, label=lst:act_cond_cpp_app]
#include "ActuatorConditioner.h"
#include <math.h>

float ActuatorConditioner::process(float rawCommand) {
    _conditionedTarget = _conditioner.process(rawCommand);
    if (!_initialized) {
        _rampedOutput = _conditionedTarget;
        _initialized = true;
    } else {
        float diff = _conditionedTarget - _rampedOutput;
        if (diff > _maxRampStep) _rampedOutput += _maxRampStep;
        else if (diff < -_maxRampStep) _rampedOutput -= _maxRampStep;
        else _rampedOutput = _conditionedTarget;
    }
    return _rampedOutput;
}
\end{lstlisting}

\subsection{MCAL Layer}

\subsubsection{Relay Driver}

\begin{lstlisting}[language=C++, caption=Relay.h --- Relay driver interface, label=lst:relay_h_app]
#ifndef RELAY_H
#define RELAY_H
#include <Arduino.h>
class Relay {
public:
    Relay(uint8_t pin, bool activeHigh = true);
    void init();
    void turnOn();
    void turnOff();
    void toggle();
    void setState(bool on);
    bool isOn() const;
private:
    uint8_t _pin; bool _activeHigh, _state;
};
#endif
\end{lstlisting}

\begin{lstlisting}[language=C++, caption=Relay.cpp --- Relay driver implementation, label=lst:relay_cpp_app]
#include "Relay.h"
Relay::Relay(uint8_t pin, bool activeHigh)
    : _pin(pin), _activeHigh(activeHigh), _state(false) {}
void Relay::init() { pinMode(_pin, OUTPUT); turnOff(); }
void Relay::turnOn() {
    _state = true;
    digitalWrite(_pin, _activeHigh ? HIGH : LOW);
}
void Relay::turnOff() {
    _state = false;
    digitalWrite(_pin, _activeHigh ? LOW : HIGH);
}
void Relay::toggle() { _state ? turnOff() : turnOn(); }
void Relay::setState(bool on) { on ? turnOn() : turnOff(); }
bool Relay::isOn() const { return _state; }
\end{lstlisting}

\subsubsection{PWM Actuator Driver}

\begin{lstlisting}[language=C++, caption=PwmActuator.h --- PWM actuator interface, label=lst:pwm_h_app]
#ifndef PWM_ACTUATOR_H
#define PWM_ACTUATOR_H
#include <Arduino.h>
class PwmActuator {
public:
    PwmActuator(uint8_t pin);
    void init();
    void setDuty(float duty);
    float getDuty() const;
    uint8_t getRawPwm() const;
private:
    uint8_t _pin; float _duty;
};
#endif
\end{lstlisting}

\begin{lstlisting}[language=C++, caption=PwmActuator.cpp --- PWM actuator implementation, label=lst:pwm_cpp_app]
#include "PwmActuator.h"
PwmActuator::PwmActuator(uint8_t pin) : _pin(pin), _duty(0.0f) {}
void PwmActuator::init() { pinMode(_pin, OUTPUT); analogWrite(_pin, 0); }
void PwmActuator::setDuty(float duty) {
    if (duty < 0.0f) duty = 0.0f;
    if (duty > 100.0f) duty = 100.0f;
    _duty = duty;
    analogWrite(_pin, (uint8_t)(_duty * 255.0f / 100.0f));
}
float PwmActuator::getDuty() const { return _duty; }
uint8_t PwmActuator::getRawPwm() const {
    return (uint8_t)(_duty * 255.0f / 100.0f);
}
\end{lstlisting}
